\chapter*{前言}

写一个操作系统很容易被误解为“先写一个内核主函数”。真正的机器没有主函数，
也不会替开发者建立栈、清空未初始化区或选择执行模式。CPU 复位以后只承诺一组
架构状态；固件必须从固定地址取到合法指令，初始化最早的观测通道，找到后续
程序，把处理器带入长模式，最后才可能进入 C++20 内核。

本书和配套工程坚持一个严格边界：QEMU 只模拟 x86-64 CPU、内存、串口、IDE
等硬件，不使用 SeaBIOS、OVMF、GRUB、Limine、Multiboot 或
\code{-kernel} 替代启动链。汇编统一采用 NASM Intel 语法，高层代码采用
freestanding C++20。构建和测试脚本运行在宿主机，可以使用 Python；进入
镜像的目标代码则必须由我们掌握。

全书既讲原理，也讲证据。每个阶段都回答四个问题：输入状态是什么，输出状态
是什么，失败怎样暴露，验收怎样自动重复。读者不必在开篇掌握全部硬件知识，
但应当愿意追踪地址、寄存器、字节布局和控制流，而不是依赖“能跑就算完成”。
